\documentclass[11pt]{article}
\usepackage[a4paper,margin=27mm]{geometry}
\usepackage{amsmath,amssymb,amsthm,booktabs,array,longtable}
\usepackage[T1]{fontenc}
\usepackage{lmodern}
\usepackage[colorlinks=true,urlcolor=blue,citecolor=blue,linkcolor=blue]{hyperref}
\newcommand{\wt}{\operatorname{wt}}
\newtheorem{theorem}{Theorem}
\newtheorem{lemma}[theorem]{Lemma}
\newtheorem{corollary}[theorem]{Corollary}
\title{Kurkov's binary Stirling row edge and extensions of
Ehrenborg--Steingr\'imsson and Erlbacher}
\author{Haobo Ma\\\small ChronoAI Pte Ltd, Singapore
\and Wenlin Zhang\\\small National University of Singapore\\\small ChronoAI Pte Ltd, Singapore}
\date{November 18, 2026}
\begin{document}
\maketitle
\begin{abstract}
We prove Kurkov's row-edge conjecture in OEIS A358612, revision 52,
FORMULA, first conjectured formula, for every nonnegative row index,
and give a polynomial refinement of Ehrenborg and Steingr\'imsson's
excedance formulas (2000, Lemma 6.2, Theorem 6.3 and Proposition 6.5)
extending the binary-product consequences in Erlbacher's note (July 2026,
Section 2 and Theorem 3.1).
The last entry is A284005$(n)/(1+\wt(n))!$, where $\wt$ is binary weight.
A dyadic block formula gives integral positive support and identifies
the complete row with a rising-factorial expansion of the polynomials
in A373183. This proves the signed factorial companion in A358612,
two Stirling-basis identities in A373183, and the conjectured identity
in MathOverflow question 483753. We also prove the lower-bit and binomial
identities in A284005 and A341392, rederive the known Boolean transform
of excedance counts, and characterize equality in a sharp lower bound.
The printed unrestricted suffix identities fail; corrected statements
hold for $0\leq q<2^m$, with a nonzero prefix in the generalized
initial-factorial version.
\end{abstract}

\section{The array and the row-edge question}
Kurkov's extended Stirling array~\cite{array} has row length
$\wt(n)+2$. Its conjectured last entry is the quotient defining
A341392~\cite{edge}. We establish both statements from the scalar
recurrence, rather than assuming the row length when computing the edge.
The block formula conjectured by Kurkov on June 21, 2024, is the
intermediate result that removes the within-row correction term.
Ehrenborg and Steingr\'imsson~\cite{excedance} give the underlying
excedance-set recurrences and inclusion--exclusion formula. Their
Proposition 6.5 supplies the all-ones special case of the signed
factorial formula. Here the polynomial identity applies to every binary
pattern and connects those counts to the literal rational array.
Erlbacher~\cite[Section 2 and Theorem 3.1]{erlbacher} derives the
binary prefix product and interprets it using ordered set partitions.
Our gap product is a regrouping of that established product, whereas
the row-edge theorem identifies it with the rational array's last entry.
Kurkov's December 2, 2020 comments~\cite{comments} already state the
equivalent product in binary run lengths.
The recurrence already appears, with its corrected plus sign, in the
November 22, 2022 MathOverflow question~\cite{precursor}.

All row and column indices are integers. Write $\mathbb N=\{0,1,\ldots\}$
and $\wt(n)$ for the number of ones in the binary expansion of $n$.
Define $T(n,k)\in\mathbb Q$ for $n,k\in\mathbb N$ by
\begin{equation}\label{initial}
 T(n,0)=0,\quad T(n,1)=1,\quad T(0,2)=1,\quad T(0,k)=0\ (k>2).
\end{equation}
For $k\geq2$, the positive-row recurrences are
\begin{align}
 T(2m+1,k)&=kT(m,k)+T(m,k-1) &&(m\geq0),\label{odd}\\
 T(2m,k)&=kT(m,k)+T(m,k-1)
 -\frac{T(2m,k-1)+T(m,k-1)}{k-1}
 &&(m\geq1).\label{even}
\end{align}
These equations determine the array uniquely: order pairs $(n,k)$
lexicographically, with the row first. In \eqref{even} the same-row term
has a smaller column; every other recursive term has a smaller row.
Column zero is a convention, not an additional source-domain equation.

Let $a(0)=1$ and
\begin{equation}\label{productrec}
 a(n)=(1+\wt(n))a(\lfloor n/2\rfloor)\qquad(n>0).
\end{equation}
This is A284005~\cite{product}, including its data value $a(1)=2$;
its NAME states the recurrence for $n>1$.
Set $e(n)=a(n)/(1+\wt(n))!$, the definition of A341392.

\begin{theorem}\label{main}
For every $n\geq0$, all $T(n,k)$ are nonnegative integers and
\[
 T(n,k)>0\ \Longleftrightarrow\ 1\leq k\leq\wt(n)+2.
 \qquad
 T(n,\wt(n)+2)=e(n).
\]
In particular, $e(n)$ is a positive integer, including $e(0)=1$.
\end{theorem}

\section{Dyadic blocks and the proof}
For $n>0$, write uniquely
\begin{equation}\label{parent}
 n=2^r(2h+1),\qquad r=v_2(n),\qquad
 h=\frac{n/2^r-1}{2}.
\end{equation}
Then $h<n$ and $\wt(n)=\wt(h)+1$. This $h$ is A025480$(n-1)$:
iterating its rule $b(2j+1)=b(j)$ exactly $r$ times gives $b(2h)=h$.

\begin{lemma}[Kurkov's block formula]\label{block}
For $r,h\geq0$ and $k\geq2$,
\begin{equation}\label{blockeq}
 T(2^r(2h+1),k)=(k-1)^rT(h,k-1)+k^{r+1}T(h,k).
\end{equation}
\end{lemma}
\begin{proof}
The case $r=0$ is \eqref{odd}. Suppose the formula holds for $r$
at every column, and put $H_k=T(h,k)$,
$U_k=T(2^r(2h+1),k)$, and $V_k=T(2^{r+1}(2h+1),k)$.
At $k=2$, \eqref{even} and $U_1=V_1=H_1=1$ give
\[
 V_2=2U_2-1=1+2^{r+2}H_2,
\]
as required. For $k\geq3$, induct on $k$, using the already established
formula for $V_{k-1}$. The two induction hypotheses give
\[
 \frac{V_{k-1}+U_{k-1}}{k-1}
 =(k-2)^rH_{k-2}+k(k-1)^rH_{k-1}.
\]
Substitute this and the formulas for $U_k,U_{k-1}$ into
$V_k=kU_k+U_{k-1}-(V_{k-1}+U_{k-1})/(k-1)$.
The $H_{k-2}$ terms and the $k(k-1)^rH_{k-1}$ terms cancel, leaving
$V_k=(k-1)^{r+1}H_{k-1}+k^{r+2}H_k$.
This proves the induction on $k$ and then on $r$.
\end{proof}

\begin{proof}[Proof of Theorem~\ref{main}]
Use strong induction on $n$. The initial row follows from
\eqref{initial} and $a(0)=1$. For $n>0$, use \eqref{parent} and put
$s=\wt(h)$. By induction the positive columns of row $h$ are precisely
$1,\ldots,s+2$ and its entries are integers. Formula~\eqref{blockeq}
is a sum of nonnegative integer multiples of these entries.
For $2\leq k\leq s+2$ its second term is positive; for $k=s+3$
its first term is positive. For $k>s+3$ both terms vanish.
Together with columns zero and one, this proves integrality and exact
support in row $n$.

At the last column the second term vanishes, so
\begin{equation}\label{edgerec}
 T(n,s+3)=(s+2)^r T(h,s+2).
\end{equation}
During the first $r+1$ halvings of $n$, the positive arguments all have
weight $s+1$, and the next argument is $h$. Thus \eqref{productrec} gives
$a(n)=(s+2)^{r+1}a(h)$. Applying the induction hypothesis to
\eqref{edgerec} now yields
\[
 T(n,s+3)=\frac{(s+2)^ra(h)}{(s+1)!}
          =\frac{a(n)}{(s+2)!}=e(n),
\]
which is the claimed edge because $\wt(n)+2=s+3$.
\end{proof}

\section{Consequences for the binary sequences}
\begin{corollary}[Binary products and binomial identities]\label{gaps}
If $n=\sum_{j=1}^{w}2^{b_j}>0$ with $0\leq b_1<\cdots<b_w$, put
$d_1=b_1$ and $d_j=b_j-b_{j-1}-1$ for $j>1$. Then
\begin{equation}\label{gapformula}
 e(n)=\prod_{j=1}^{w}(w-j+2)^{d_j}.
\end{equation}
The empty product gives $e(0)=1$. Furthermore, for $m,n\geq0$,
\begin{align}
 e(2n+1)&=e(n),&e(2n)&=(1+\wt(n))e(n),\label{erecs}\\
 e(2^m(2n+1))&=\sum_{j=0}^{m}\binom mj e(2^j n),\label{ebin}\\
 a(2^m(2n+1))&=\sum_{j=0}^{m+1}\binom{m+1}j a(2^j n).\label{abin}
\end{align}
\end{corollary}
\begin{proof}
Equation~\eqref{edgerec} gives $e(2^r(2h+1))=(\wt(h)+2)^re(h)$.
Successively remove the least significant one and its following right
shift. The exponents encountered are $d_1,\ldots,d_w$ and the bases
are $w+1,w,\ldots,2$, proving \eqref{gapformula}.
The odd recurrence in \eqref{erecs} follows by taking $r=0$;
the even recurrence follows from \eqref{productrec}, with $n=0$
checked separately.

Consequently $e(2^j n)=(1+\wt(n))^je(n)$, also for $n=0$.
The binomial theorem makes the right side of \eqref{ebin}
$(2+\wt(n))^me(n)$, which equals its left side by the dyadic edge
recurrence. Similarly,
$a(2^j n)=(1+\wt(n))^ja(n)$, including $n=0$, and
$a(2^m(2n+1))=(2+\wt(n))^{m+1}a(n)$.
The binomial theorem proves \eqref{abin}. These are the April 24, 2023
displayed identities in A341392 and A284005, respectively.
The product \eqref{gapformula} is the factorial normalization of
the prefix product derived by Erlbacher~\cite[Section 2, (3)--(4)]{erlbacher}.
\end{proof}

The same one-step formulas give $a(2n+1)=a(n)+a(2n)$ for $n\geq0$.
They also give the displayed even-step identities in both entries:
\begin{equation}\label{lower}
 a(2n)=a(n)+a(2n-2^{v_2(n)}),\qquad
 e(2n)=e(n)+e(2n-2^{v_2(n)})\quad(n>0).
\end{equation}
Indeed, write $n=2^r(2h+1)$ and $s=\wt(h)$, so that
$2n-2^r=2^r(4h+1)$. The dyadic recurrences give
\[
\begin{aligned}
 a(n)&=(s+2)^{r+1}a(h),&
 a(2n-2^r)&=(s+2)^{r+1}(s+1)a(h),\\
 e(n)&=(s+2)^re(h),&
 e(2n-2^r)&=(s+2)^r(s+1)e(h).
\end{aligned}
\]
Thus each second expression is $\wt(n)$ times the first, proving
\eqref{lower} by the doubling recurrences. These statements include the
case $h=0$ because $a(0)=e(0)=1$ and $a(2h)=(1+s)a(h)$,
$e(2h)=(1+s)e(h)$ also hold at zero.

For example, $10=2^1+2^3$ gives $e(10)=3\cdot2=6$, and row ten is
$(1,21,28,6)$. The two remaining elementary formulas in A358612 also
follow: $T(n,2)=2n+1$ by induction using \eqref{odd} and
$T(2m,2)=2T(m,2)-1$; and
$T(2^j-1,k)=S(j+2,k)$ for $j\geq0,k\geq1$, by repeatedly using
\eqref{odd} and the Stirling recurrence~\cite[Eq.~26.8.22]{dlmf},
starting with $S(2,1)=S(2,2)=1$.

\begin{corollary}[Lower-bit identities]\label{bitflip}
Let $n=2^L+h>1$, $0\leq h<2^L$, and write $\epsilon_j$ for bit $j$
of $h$. Then
\begin{align}
 a(n)&=2a(h)+\sum_{j=0}^{L-1}a\bigl(h+2^j(1-\epsilon_j)\bigr),\label{aflip}\\
 e(n)&=e(h)+\sum_{\substack{0\leq j<L\\\epsilon_j=0}}e(h+2^j).
 \label{eflip}
\end{align}
These are Kurkov's November 10, 2019 conjecture in A284005 and the
April 24, 2023 formula in A341392, respectively. The former also answers
the proof request in MathOverflow question 445013~\cite{lowerquestion}.
\end{corollary}
\begin{proof}
Put $c_j=1+\wt(\lfloor h/2^j\rfloor)$, so $c_L=1$ and
$c_{j+1}=c_j-\epsilon_j$. Let $p_0=1$ and
$p_{j+1}=p_j(c_j+1)/c_j$. The product recurrence gives
$a(n)=2p_La(h)$ and, when $\epsilon_j=0$,
$a(h+2^j)=p_{j+1}a(h)$. Since
\[
 (c_{j+1}+1)p_{j+1}-(c_j+1)p_j=(1-\epsilon_j)p_{j+1},
\]
summing yields $2p_L=\wt(h)+2+\sum_{\epsilon_j=0}p_{j+1}$.
This proves \eqref{aflip}, since its one-bit summands equal $a(h)$.
Divide the equivalent identity
$a(n)=(\wt(h)+2)a(h)+\sum_{\epsilon_j=0}a(h+2^j)$ by
$(\wt(h)+2)!$ to obtain \eqref{eflip}. The same computation covers $h=0$.
\end{proof}

\section{The complete polynomial identity}
For a polynomial $f\in\mathbb Q[x]$, define
$Xf(x)=xf(x)$ and $Df(x)=x(f(x+1)-f(x))$. Define the polynomials
of A373183~\cite{poly} by
\[
 R_0(x)=x,\qquad R_{2h+1}=XR_h,\qquad R_{2h}=DR_h\quad(h>0).
\]
The even rule also holds at $h=0$, because $Dx=x$.
Write $P(n,q)=[x^q]R_n(x)$ for the coefficient of $x^q$.
For $i\geq1$, let $x^{\overline i}=x(x+1)\cdots(x+i-1)$,
and let $(x)_i=x(x-1)\cdots(x-i+1)$.
The signed Stirling numbers $s(i,q)$ are defined by
$(x)_i=\sum_q s(i,q)x^q$.

\begin{theorem}[Polynomial refinement]\label{polynomial}
For every $n\geq0$, putting $w=\wt(n)$, we have
\begin{equation}\label{rising}
 xR_n(x)=\sum_{i=1}^{w+2}(-1)^{w+2-i}T(n,i)x^{\overline i}.
\end{equation}
In particular, for $r,h\geq0$,
\begin{equation}\label{risingblock}
 R_{2^r(2h+1)}(x)=\sum_{i=1}^{\wt(h)+2}
 (-1)^{\wt(h)+2-i}T(h,i)i^r x^{\overline i}.
\end{equation}
\end{theorem}
\begin{proof}
The elementary identities
\[
 D(x^{\overline i})=i x^{\overline i},\qquad
 x^{\overline{i+1}}=(x+i)x^{\overline i}
\]
follow directly from the products. Denote the right side of
\eqref{rising} by $G_n(x)$. Its initial value is
$G_0=x^{\overline2}-x^{\overline1}=x^2=xR_0$.
For $n=2^r(2h+1)>0$, use Lemma~\ref{block} and
$\wt(n)=\wt(h)+1$. Reindex the term involving $T(h,i-1)$.
The two resulting terms for index $i$ combine to give
\[
 G_n=\sum_i(-1)^{\wt(h)+2-i}T(h,i)i^r
       \bigl(x^{\overline{i+1}}-i x^{\overline i}\bigr)
     =xD^rG_h.
\]
There is no boundary term: $T(h,0)=0$, and columns above the support
vanish. The column-one term is $T(n,1)=T(h,1)=1$.
Strong induction on $n$ gives $G_h=xR_h$, and hence
$G_n=xD^r(XR_h)=xR_n$. Applying the eigenvalue identity to $G_h$
also proves \eqref{risingblock}.
\end{proof}

\begin{corollary}[The two Stirling-basis conjectures]\label{basis}
For all $n\geq0$ and $q>0$,
\begin{equation}\label{coefficient}
 P(n,q)=(-1)^{\wt(n)-q+1}
       \sum_{i=q+1}^{\wt(n)+2}s(i,q+1)T(n,i).
\end{equation}
For $r,h\geq0$ and $q>0$,
\begin{equation}\label{coefficientblock}
 P(2^r(2h+1),q)=(-1)^{\wt(h)-q}
       \sum_{i=q}^{\wt(h)+2}s(i,q)T(h,i)i^r.
\end{equation}
These prove, respectively, the final June 19, 2024 formula and the
January 17, 2025 conjecture in A373183, revision 28, FORMULA.
Moreover,
\begin{equation}\label{falling}
 \frac1x\sum_{i=1}^{\wt(n)+2}T(n,i)(x)_i
 =(-1)^{\wt(n)+1}R_n(-x).
\end{equation}
Thus the conjectured identity in MathOverflow question 483753,
December 7, 2024~\cite{mo}, holds as a polynomial identity.
\end{corollary}
\begin{proof}
The coefficient of $x^q$ in $x^{\overline i}$ is
$(-1)^{i-q}s(i,q)$. Extract coefficients in \eqref{rising} and
\eqref{risingblock}. Replacing $x$ by $-x$ in \eqref{rising} and
using $(-x)^{\overline i}=(-1)^i(x)_i$ gives \eqref{falling}.
Each falling factorial is divisible by $x$, so the quotient is a
polynomial, including at $x=0$.
\end{proof}

\section{Excedance counts and Boolean transforms}
Let $b(n)$ denote A329369~\cite{excentry}, the excedance-set count
specified by the binary digits of $n$. Its established recurrences are
$b(0)=1$, $b(2h+1)=b(h)$ and, for $h>0$,
\begin{equation}\label{brec}
 b(2h)=b(h)+b(h-2^{v_2(h)})+b(2h-2^{v_2(h)}).
\end{equation}
They are the word recurrences of Ehrenborg and
Steingr\'imsson~\cite[Propositions 2.1 and 2.5]{excedance}.

\begin{corollary}[Signed factorial companion]\label{signed}
For $r,h\geq0$,
\begin{equation}\label{factorialsum}
 b(2^r(2h+1))=\sum_{i=1}^{\wt(h)+2}
 (-1)^{\wt(h)+2-i}i!i^rT(h,i).
\end{equation}
This proves the second conjectured formula in A358612, revision 52,
FORMULA. For $h=2^j-1$, it specializes to Ehrenborg and
Steingr\'imsson's formula for two-block words~\cite[Proposition 6.5]{excedance}.
\end{corollary}
\begin{proof}
Direct expansion gives $DX=X+ D+XD$. Write $h=2^t(2u+1)>0$.
Apply $D^t$ to this operator identity acting on $R_u$. The result is
\[
 R_{2h}=R_h+R_{h-2^t}+R_{2h-2^t},
\]
because $h-2^t=2^{t+1}u$ and $2h-2^t=2^t(4u+1)$.
This also holds at $u=0$, using $D^jR_0=R_0$.
The numbers $R_n(1)$ therefore satisfy \eqref{brec}, the odd rule,
and the same initial value as $b(n)$. These rules uniquely determine
$b$ by induction on the index, so $R_n(1)=b(n)$.
Evaluate \eqref{risingblock} at $x=1$, where
$1^{\overline i}=i!$. The all-ones specialization uses
$T(2^j-1,i)=S(j+2,i)$.
\end{proof}

Write $k\subseteq_2 n$ when every one-bit of $k$ is a one-bit of $n$.
\begin{corollary}[Boolean transform; known excedance consequence]\label{boolean}
For $n\geq0$,
\begin{equation}\label{bool}
 a(n)=\sum_{k\subseteq_2 n}b(k),\qquad
 b(n)=\sum_{k\subseteq_2 n}(-1)^{\wt(n)-\wt(k)}a(k).
\end{equation}
The first identity is Kurkov's March 10, 2023 conjecture in A284005;
the second is the inverse-transform conjecture in A329369.
They are consequences of Ehrenborg and Steingr\'imsson's
Lemma 6.2 and Theorem 6.3 (2000); the following proof gives their
operator form in the present notation.
JetfiRex's later answer~\cite{jetfirex} discusses the same transform
through constrained permutations and Boolean inversion.
\end{corollary}
\begin{proof}
Read the binary digits from most to least significant, applying $X$
for a one and $D$ for a zero to the initial polynomial $x$.
This constructs $R_n$. Summing over all submasks replaces each $X$
by $D+X$. On the rising-factorial basis,
$(D+X)x^{\overline i}=x^{\overline{i+1}}$ and
$Dx^{\overline i}=i x^{\overline i}$.
Consequently the sum evaluates at $x=1$ to
$(\wt(n)+1)!\prod_{j=1}^{\wt(n)}(\wt(n)-j+2)^{d_j}=a(n)$,
by Corollary~\ref{gaps}. Leading zero padding acts as $Dx=x$.
Boolean inclusion--exclusion gives the inverse identity, since for
$\ell\subseteq_2 n$ the sum
$\sum_{\ell\subseteq_2 k\subseteq_2 n}(-1)^{\wt(n)-\wt(k)}$
equals $[\ell=n]$. Finally, Lucas's parity rule identifies
$k\subseteq_2 n$ with $\binom nk$ odd: over $\mathbb F_2$,
$(1+z)^n=\prod_{j:\,n_j=1}(1+z^{2^j})$.
\end{proof}

\begin{corollary}[General initial factorials]\label{initialfactorials}
Let $\mathcal L_n$ be the operator word that constructs $R_n$ from $x$,
reading the shortest binary expansion of $n$, with $\mathcal L_0=I$.
For an integer $k>0$ set $U_{n,k}=\mathcal L_n x^{\overline k}$.
Then A357990$(n,k)=U_{n,k}(1)/k!$, and in particular
A357990$(n,1)=b(n)$, its conjectured first column~\cite{square}.
If
\[
 d(n,k,j)=[x^{\wt(n)+k-j+1}]U_{n,k}(x),
\]
where coefficients outside the polynomial are zero, these are precisely
Kurkov's $d$ coefficients in the June 19, 2024 formulas of A373183.
Thus those formulas for $P(n,q)$, for the generalized initial polynomial,
and for $d(n,1,1)=e(n)$ hold.
\end{corollary}
\begin{proof}
For $n=2^L+h>0$, let $t$ be the bit length of $h$, with $t=0$ at zero,
and put $g=L-t$. The leading-one decomposition is
$\mathcal L_n=\mathcal L_hD^gX$. Therefore
\[
 U_{n,k}=(k+1)^gU_{h,k+1}-k^{g+1}U_{h,k}.
\]
After evaluation at one and division by $k!$, this becomes
$(k+1)^{g+1}F_h(k+1)-k^{g+1}F_h(k)$, with $F_0(k)=1$.
This is exactly the defining recurrence of A357990.
Taking coefficients in reverse degree gives
\[
 d(n,k,j)=(k+1)^gd(h,k+1,j)-k^{g+1}d(h,k,j-1).
\]
At zero, $d(0,k,j)=|s(k,k-j+1)|$ for $1\leq j\leq k$, and
$d(n,k,0)=0$. These agree with the source's initial values and recurrence.
For $k=1$, $U_{n,1}=R_n$, whose leading coefficient is $e(n)$ as shown
in Corollary~\ref{sharp} below.
\end{proof}

\begin{corollary}[A corrected suffix range]\label{suffix}
The second conjectured formula in A373183, revision 28, and
Conjecture 6 in A329369, revision 236 (June 5, 2024), hold with the range
\begin{equation}\label{suffixrange}
 n,m\geq0,\qquad 0\leq q<2^m:
 \qquad b(2^mn+q)=\sum_{i=1}^{\wt(n)+1}
 P(n,i)b(2^m(2^{i-1}-1)+q).
\end{equation}
Their printed range $q\geq0$, with no upper bound, fails at
$(n,m,q)=(2,0,1)$: the left side is $1$ and the right side is $7$.
For the June 19, 2024 generalized $d$ sum in A373183, a valid range is
$n>0$, $k>0$, $m\geq0$, $0\leq q<2^m$, and the identity is
\begin{equation}\label{dsuffix}
 k!\,\mathrm{A357990}(2^mn+q,k)
 =\sum_{i=1}^{\wt(n)+k}d(n,k,\wt(n)-i+k+1)
       b(2^m(2^{i-1}-1)+q).
\end{equation}
The printed generalized range also fails at $(n,k,m,q)=(0,2,1,0)$:
the two sides are $2$ and $4$ respectively.
\end{corollary}
\begin{proof}
Let $\mathcal S_{m,q}$ be the operator word for the $m$-digit binary
suffix $q$, including leading zeros. For $q<2^m$, concatenation gives
$R_{2^mn+q}=\mathcal S_{m,q}R_n$.
This remains valid at $n=0$ because $Dx=x$.
Since $R_{2^{i-1}-1}=x^i$, applying the linear functional
$f\mapsto(\mathcal S_{m,q}f)(1)$ to
$R_n=\sum_iP(n,i)x^i$ proves \eqref{suffixrange}.
For $n>0$, the same concatenation gives
$U_{2^mn+q,k}=\mathcal S_{m,q}U_{n,k}$ and proves \eqref{dsuffix}.
The source's lower limit $v_2(n+1)+1$ can be restored because all lower
coefficients vanish: an odd digit shifts the minimum degree, and a zero
digit resets it to one.
For the first counterexample, $R_2=x+2x^2$ and
$(b(0),b(1),b(2),b(3))=(1,1,3,1)$, so $1\ne1+2\cdot3$.
For the second, $U_{0,2}=x+x^2$, so the sum is $b(0)+b(2)=4$,
whereas $2!\,\mathrm{A357990}(0,2)=2$.
At $n=0$ the generalized concatenation cannot discard leading zeros:
$Dx^{\overline k}=kx^{\overline k}$, rather than $x^{\overline k}$.
\end{proof}

\begin{corollary}[Sharp comparison with the edge]\label{sharp}
For all $n\geq0$,
\[
 b(n)\geq2e(n)-1.
\]
For $n>0$, equality holds exactly when deleting the leading one-bit
leaves only trailing ones (zero is allowed). Equality also holds at zero.
\end{corollary}
\begin{proof}
The coefficients of $R_n$ are positive exactly in degrees
$v_2(n+1)+1,\ldots,\wt(n)+1$. Indeed $X$ shifts the support by one;
for a polynomial with nonnegative coefficients and positive leading
coefficient in degree $d>0$, $D$ has positive coefficients in every
degree $1,\ldots,d$, since
$Dx^j=\sum_{q=1}^j\binom j{q-1}x^q$.
This proves the support statement by induction, starting with $R_0=x$.
The leading coefficient is $e(n)$: $X$ preserves it, $D$ multiplies it
by the degree, and these are exactly \eqref{erecs}.
Also $R_n(-1)=(-1)^{\wt(n)+1}$, by the odd rule and
$DR_n(-1)=R_n(-1)$, since $R_n(0)=0$.
Writing $d=\wt(n)+1$ therefore gives
\[
 b(n)=2e(n)-1+
 2\sum_{\substack{1\leq q<d\\q\equiv d\pmod2}}P(n,q).
\]
The sum vanishes exactly when $v_2(n+1)+1>d-2$, or
$v_2(n+1)\geq\wt(n)-1$. For $n>0$ this is precisely the stated
binary condition. At $n=0$ both sides equal one.
\end{proof}
Since A036987 is the indicator of integers $t$ for which $t+1$ is a
power of two~\cite{indicator}, this criterion proves the bracketed
comparison in A341392: the indicator of $2e(n)-1=b(n)$ is
A036987$(\mathrm{A053645}(n))$, where A053645 deletes the leading
one-bit and maps zero to zero. The unrestricted equality
$2e(n)-1=b(n)$ is not asserted by that bracketed formula.

\begin{corollary}[Sums by binary weight]\label{eulerian}
For $L,k\geq0$,
\[
 \sum_{\substack{0\leq n<2^L\\\wt(n)=k}}b(n)
 =\genfrac{\langle}{\rangle}{0pt}{}{L+1}{k},
 \qquad \sum_{0\leq n<2^L}b(n)=(L+1)!,
\]
where the Eulerian number in brackets counts permutations of $L+1$
elements with exactly $k$ excedances.
\end{corollary}
\begin{proof}
The $L$ binary digits specify each subset of $\{1,\ldots,L\}$ exactly
once, and $b(n)$ counts the permutations having that excedance set
\cite{excedance,excentry}. Grouping by cardinality proves the first
identity and summing over all subsets proves the second. Thus
Conjecture 7 in A329369, dated January 27, 2025, follows directly from
the established excedance interpretation.
\end{proof}

\section{Coverage and verification}
The source items covered by the argument are distinguished below from
related questions that require a different argument. ``Not covered''
specifies this note's scope, not the current status of those questions.

\begingroup\small
\begin{longtable}{>{\raggedright\arraybackslash}p{0.56\linewidth}>{\raggedright\arraybackslash}p{0.36\linewidth}}
\toprule
Source item & Coverage\\\midrule
A358612, row length and first conjectured formula & Proved, Theorem~\ref{main}\\
A358612, June 21, 2024 block conjecture & Proved, Lemma~\ref{block}\\
A358612, column two and all-ones rows & Proved above\\
Binary-gap product for A341392 & Reproved; prior product proof Erlbacher (2026), Section 2\\
A341392, April 24, 2023 binomial identity & Proved, \eqref{ebin}\\
A284005, April 24, 2023 binomial conjecture & Proved, \eqref{abin}\\
A341392/A284005, displayed odd and even-step recurrences & Proved, \eqref{erecs}, \eqref{lower} and the text\\
A358612, signed factorial sum; A329369 Conjecture 4 & Proved, Corollary~\ref{signed}; two-block case reproved, prior proof Ehrenborg--Steingr\'imsson (2000), Proposition 6.5\\
A284005, November 10, 2019 lower-bit conjecture; A341392, April 24, 2023 bit-flip formula & Proved, Corollary~\ref{bitflip}\\
A284005, March 10, 2023 Boolean transform; A329369 inverse-transform conjecture & Reproved, Corollary~\ref{boolean}; prior proof Ehrenborg--Steingr\'imsson (2000), Lemma 6.2 and Theorem 6.3\\
A373183, June 19, 2024 Stirling-basis formula and January 17, 2025 conjecture & Proved, Corollary~\ref{basis}\\
MathOverflow 483753, displayed conjecture and request for a proof & Proved; answered: yes, Corollary~\ref{basis}\\
A357990 first-column conjecture; A329369 Conjecture 3 & Proved, Corollary~\ref{initialfactorials}\\
A373183, June 19, 2024 $d$ recurrence, initial-factorial extension and leading coefficient & Proved, Corollary~\ref{initialfactorials}\\
A373183, first three-term polynomial conjecture & Proved in Corollary~\ref{signed}\\
A373183 second conjecture; A329369 Conjecture 6 (unrestricted suffix) & Refuted-and-corrected, Corollary~\ref{suffix}, \eqref{suffixrange}\\
A373183, June 19, 2024 generalized $d$ sum & Refuted-and-corrected, Corollary~\ref{suffix}, \eqref{dsuffix}\\
A341392, bracketed equality comparison with A329369 & Proved by the equality criterion in Corollary~\ref{sharp}\\
The bound $b(n)\geq2e(n)-1$ & Not posed; proved, Corollary~\ref{sharp}\\
A329369, factorial row sum and Conjecture 7 & Reproved; prior count interpretation Ehrenborg--Steingr\'imsson (2000), Corollary~\ref{eulerian}\\
MathOverflow 445013, request for a lower-bit recurrence proof & Answered: yes, Corollary~\ref{bitflip}\\
A329369, July 10, 2024 two-parameter Stirling recursion & Prior proof cited: Alekseyev (2024)~\cite{alekseyev}; not needed here\\
A373183, long $W$-coefficient expansion; A329369 rook-tour interpretation and Conjectures 5, 8--10 & Not covered\\
A341392, A243499 value-set/coordinate identities and FindStat mappings; A284005 divisor-count identity & Not covered\\
A358631 first-kind array & Not covered\\
MathOverflow 486535, sums of the binary Stirling coefficients~\cite{aggregate} & Not covered\\
\bottomrule
\end{longtable}
\endgroup

\begingroup\raggedright
Every result claimed here has a written proof. The pinned Lean source
\href{https://github.com/the-omega-institute/trureturing/blob/47f05ad67b37a0f7cbc4673c33bb17fb7c1e78fb/D5/S1/Recurrence/Parity/DyadicExtendedStirlingEdge.lean}{\texttt{DyadicExtendedStirlingEdge.lean}}
contains \texttt{support}, the vanishing statement above the binary-weight
edge, and \texttt{result}, the rational edge identity for all $n\in\mathbb N$.
The definition of $T$ and declaration \texttt{dyadic\_extended\_stirling\_block}
are in the pinned
\href{https://github.com/the-omega-institute/trureturing/blob/47f05ad67b37a0f7cbc4673c33bb17fb7c1e78fb/D5/S1/Recurrence/Parity/DyadicExtendedStirlingBlock.lean}{\texttt{DyadicExtendedStirlingBlock.lean}}.
The definition of $a$, its dyadic identity \texttt{a\_two\_pow\_mul\_odd},
and \eqref{abin} as \texttt{result} are in the pinned
\href{https://github.com/the-omega-institute/trureturing/blob/47f05ad67b37a0f7cbc4673c33bb17fb7c1e78fb/D5/S1/Digit/KurkovWeightProductBinomial.lean}{\texttt{KurkovWeightProductBinomial.lean}}.
Here \texttt{wt} is the binary digit sum defined in the pinned
\href{https://github.com/the-omega-institute/trureturing/blob/47f05ad67b37a0f7cbc4673c33bb17fb7c1e78fb/D5/S1/Digit/DyadicRowPolynomialRecurrence.lean}{\texttt{DyadicRowPolynomialRecurrence.lean}}.
That file also contains \texttt{D\_X\_mul}, the operator relation
$DX=X+(X+1)D$, and \texttt{conjecture1}, the three-term polynomial
recurrence used in Corollary~\ref{signed}. Its \texttt{R\_degree}
is an upper bound, not the exact positive-support statement.
Strict positivity, integrality, the gap and lower-bit formulas, the complete
rising-factorial connection, all coefficient and signed-transform
consequences, the corrected suffix ranges and the sharp inequality have
the written proofs given here. The edge file's \texttt{support} asserts
vanishing above the edge, not positivity below it.\par
\endgroup

\begin{samepage}
The accompanying \texttt{verify.py} evaluates the original recurrence in
exact rational arithmetic for $0\leq n\leq4095$, $0\leq k\leq17$.
It checks the block formula, exact support, edge identity, polynomial
identities, signed and Boolean transforms, lower-bit formulas, corrected
suffix identities, counterexamples and the sharp bound on its stated
finite ranges. The output \texttt{verification.json} records those ranges;
these checks supplement the all-index proofs.\par
\end{samepage}

\begin{thebibliography}{99}
\bibitem{array} M. Kurkov, \emph{Irregular triangle of extended Stirling
numbers of the second kind}, OEIS A358612, revision 52, March 16, 2026;
FORMULA and COMMENTS.
\href{https://raw.githubusercontent.com/oeis/oeisdata/d8ad22d9b98d2c427be8c6cd9f2672c600b8e3ae/seq/A358/A358612.seq}{Version-specific entry}.
\bibitem{edge} M. Kurkov, \emph{A284005$(n)/(1+\mathrm{A000120}(n))!$},
OEIS A341392, revision 56; NAME and FORMULA.
\href{https://raw.githubusercontent.com/oeis/oeisdata/d8ad22d9b98d2c427be8c6cd9f2672c600b8e3ae/seq/A341/A341392.seq}{Version-specific entry}.
\bibitem{product} A. Karttunen, \emph{Binary-weight product sequence},
OEIS A284005, revision 106; NAME, DATA, and Kurkov's FORMULA of
April 24, 2023.
\href{https://raw.githubusercontent.com/oeis/oeisdata/d8ad22d9b98d2c427be8c6cd9f2672c600b8e3ae/seq/A284/A284005.seq}{Version-specific entry}.
\bibitem{dlmf} NIST, \emph{Digital Library of Mathematical Functions},
\S26.8, Set Partitions: Stirling Numbers, Eq.~26.8.22.
\url{https://dlmf.nist.gov/26.8.E22}.
\bibitem{mo} user231922, \emph{Generalized identity with Stirling numbers
of the second kind and falling factorials}, MathOverflow question 483753,
December 7, 2024. \url{https://mathoverflow.net/q/483753}.
\bibitem{poly} M. Kurkov, \emph{Irregular triangle with row polynomials},
OEIS A373183, revision 28, March 16, 2026; NAME and FORMULA.
\href{https://raw.githubusercontent.com/oeis/oeisdata/d8ad22d9b98d2c427be8c6cd9f2672c600b8e3ae/seq/A373/A373183.seq}{Version-specific entry}.
\bibitem{excentry} M. Kurkov, \emph{Number of permutations with a prescribed
binary excedance set}, OEIS A329369, revision 236, March 14, 2026;
NAME, COMMENTS and FORMULA.
\href{https://raw.githubusercontent.com/oeis/oeisdata/d8ad22d9b98d2c427be8c6cd9f2672c600b8e3ae/seq/A329/A329369.seq}{Version-specific entry}.
\bibitem{excedance} R. Ehrenborg and E. Steingr\'imsson,
\emph{The excedance set of a permutation}, Advances in Applied Mathematics
\textbf{24} (2000), 284--299. \href{https://doi.org/10.1006/aama.1999.0671}{doi:10.1006/aama.1999.0671}.
\href{https://www.ms.uky.edu/~jrge/papers/Excedance.pdf}{Author-hosted full text}.
\bibitem{square} M. Kurkov, \emph{Square array defined by a leading-bit
recurrence}, OEIS A357990, revision 41, March 29, 2026; NAME and FORMULA.
\href{https://raw.githubusercontent.com/oeis/oeisdata/d8ad22d9b98d2c427be8c6cd9f2672c600b8e3ae/seq/A357/A357990.seq}{Version-specific entry}.
\bibitem{alekseyev} M. Alekseyev, \emph{Recursion for the sum with Stirling
numbers of both kinds}, MathOverflow answer 474668, July 8, 2024.
\url{https://mathoverflow.net/a/474668}.
\bibitem{indicator} \emph{Fredholm--Rueppel sequence}, OEIS A036987,
revision 243, November 28, 2025; COMMENTS and FORMULA.
\href{https://raw.githubusercontent.com/oeis/oeisdata/d8ad22d9b98d2c427be8c6cd9f2672c600b8e3ae/seq/A036/A036987.seq}{Version-specific entry}.
\bibitem{erlbacher} J. Erlbacher, \emph{A proof of Kurkov's conjecture for
A000670}, July 2026, Section 2 and Theorem 3.1.
\href{https://raw.githubusercontent.com/demonstrandum-research/artifacts/de462739d3afc6bfb3316d69a840ab4f75007aa9/papers/kurkov-a000670/note.pdf}{Version-specific full text}.
\bibitem{precursor} user231922, \emph{Closed form for coefficients related
to excedance set of permutation}, MathOverflow question 435104,
November 22, 2022; recurrence corrected that day at 11:15:51 UTC.
\url{https://mathoverflow.net/q/435104}.
\bibitem{comments} M. Kurkov, \emph{Comments on A284005},
December 2, 2020; binary run-length product.
\url{https://oeis.org/A284005/a284005.txt}.
\bibitem{jetfirex} JetfiRex, answer to \emph{Sum with products turned
into subsequences}, MathOverflow answer 413850, January 14, 2022,
edited January 18, 2022. \url{https://mathoverflow.net/a/413850}.
\bibitem{lowerquestion} user231922, \emph{Recurrence for the A284005},
MathOverflow question 445013, April 18, 2023.
\url{https://mathoverflow.net/q/445013}.
\bibitem{aggregate} user231922, \emph{Integer coefficients that sum up to
Stirling transform of shifted A001475}, MathOverflow question 486535,
January 24, 2025. \url{https://mathoverflow.net/q/486535}.
\end{thebibliography}
\noindent\small\copyright\ 2026 Haobo Ma and Wenlin Zhang.
Licensed under \href{https://creativecommons.org/licenses/by/4.0/}{CC BY 4.0}.
\end{document}
